\chapter{Himpunan dan Struktur}\label{ch:b2:structures}

Bab pembuka ini mempertajam landasan yang diletakkan pada jilid Tahun
ke-1 menjadi perkakas kerja sehari-hari: kalkulus himpunan dan kuosien,
perbandingan himpunan tak hingga (keterbilangan, Cantor--Bernstein),
serta teori struktur grup dan ring --- teorema Lagrange, grup simetri
beserta tanda permutasinya, \omterm{def:b2:structures:ideal}{ideal}, dan teorema sisa Cina. Semua yang ada
di sini dipakai tiada henti pada sisa buku: tanda permutasi membangun
determinan (\cref{ch:b2:linalg}), \omterm{def:b2:structures:quotientring}{ring kuosien} menggerakkan aritmetika,
dan keterbilangan menjadi alas bagi topologi maupun peluang.

\section{Himpunan, pemetaan, kuosien}

Kita memakai dengan bebas bahasa himpunan, pemetaan, serta relasi
ekuivalensi dan relasi urutan yang disiapkan pada jilid Tahun ke-1. Dua
peningkatan layak dinyatakan secara utuh.

\begin{proposition}[Peta dan prapeta suatu keluarga]\label{prop:b2:structures:images}
Misalkan $f \colon E \to F$ dan misalkan $(A_i)_{i \in I}$, $(B_j)_{j \in J}$
keluarga himpunan bagian dari $E$ dan dari $F$ berturut-turut. Maka
\[
f^{-1}\Bigl(\bigcup_j B_j\Bigr) = \bigcup_j f^{-1}(B_j),
\qquad
f^{-1}\Bigl(\bigcap_j B_j\Bigr) = \bigcap_j f^{-1}(B_j),
\qquad
f^{-1}(F \setminus B) = E \setminus f^{-1}(B),
\]
\[
f\Bigl(\bigcup_i A_i\Bigr) = \bigcup_i f(A_i),
\qquad
f\Bigl(\bigcap_i A_i\Bigr) \subseteq \bigcap_i f(A_i)
\quad (\text{kesamaan bila } f \text{ injektif}).
\]
\end{proposition}

\begin{proof}
Setiap kesamaan hanyalah penguraian definisi; misalnya $x \in
f^{-1}(\bigcap B_j) \iff f(x) \in B_j$ untuk setiap $j$ $\iff x \in
f^{-1}(B_j)$ untuk setiap $j$. Kesamaan untuk peta dan gagalnya
kesamaan pada kasus irisan (beserta perbaikan lewat keinjektifan) telah
dibuktikan pada jilid Tahun ke-1 untuk dua himpunan; alasannya sama
persis untuk keluarga sembarang.
\end{proof}

\begin{example}[Ketika pengaitan peta benar-benar tegas]\label{ex:b2:structures:imagestrict}
Ambil $f \colon \R \to \R$, $f(x) = x^2$, dengan $A_1 =
\intcc{-1}{0}$ dan $A_2 = \intcc{0}{1}$. Maka
\[
f(A_1 \cap A_2) = f(\{0\}) = \{0\},
\qquad
f(A_1) \cap f(A_2) = \intcc{0}{1} \cap \intcc{0}{1} =
\intcc{0}{1} :
\]
pengaitan pada \cref{prop:b2:structures:images} setegas yang
mungkin --- kedua prapeta $\pm x$ dari satu nilai yang sama
berada di $A_i$ yang berbeda. Keinjektifan justru melarang
pemecahan semacam ini, dan itulah sebabnya prapeta (yang tak
pernah menyatukan titik) memenuhi keempat kesamaan itu tanpa syarat,
sedangkan peta kehilangan kesamaan untuk irisan. Pedoman praktis
untuk seluruh buku: dorong \emph{prapeta} melewati operasi
himpunan sesuka hati; perlakukan peta dengan hati-hati.
\end{example}

\begin{definition}[Himpunan kuosien]\label{def:b2:structures:quotient}
Misalkan $\mathcal{R}$ relasi ekuivalensi pada $E$. \emph{Himpunan
kuosien}\index{himpunan kuosien} $E/\mathcal{R}$ adalah himpunan semua
kelas ekuivalensi; surjeksi $\pi \colon E \to E/\mathcal{R}$,
$x \mapsto \mathrm{cl}(x)$, disebut \emph{proyeksi
kanonik}\index{proyeksi kanonik}.

\emph{Sifat universal (pemfaktoran):} jika $f \colon E \to F$
\emph{selaras} dengan $\mathcal{R}$ (yaitu $x \mathbin{\mathcal{R}}
y \implies f(x) = f(y)$), maka terdapat tepat satu pemetaan $\overline f
\colon E/\mathcal{R} \to F$ dengan $f = \overline f \circ \pi$.
\end{definition}

\begin{proof}[Bukti sifat universal]
Ketunggalan: syarat $f = \overline f \circ \pi$ berbunyi
\[
\overline f\bigl(\mathrm{cl}(x)\bigr) = f(x)
\qquad (x \in E),
\]
dan karena $\pi$ surjektif, setiap unsur
$E/\mathcal{R}$ berbentuk $\mathrm{cl}(x)$: seluruh nilai
$\overline f$ sudah terpaksa. Keberadaan: ambil ungkapan di atas
sebagai \emph{definisi} $\overline f$; ungkapan itu tak bermakna
ganda justru karena keselarasan --- jika $\mathrm{cl}(x) =
\mathrm{cl}(y)$, maka $x \mathbin{\mathcal{R}} y$, sehingga $f(x) =
f(y)$ dan kedua calon nilai berimpit --- dan ungkapan itu
memfaktorkan $f$ menurut konstruksinya. Perhatikan pembagian tugas:
kesurjektifan $\pi$ memberi ketunggalan, keselarasan memberi keberadaan.
\end{proof}

\begin{example}
$\Z/n\Z$ adalah kuosien $\Z$ oleh kekongruenan modulo $n$; pemeriksaan
``terdefinisi dengan baik'' pada jilid Tahun ke-1 tak lain adalah
penerapan sifat universal. Kuosien mengubah ``konstruksi yang selaras
pada wakil kelas'' menjadi pemetaan sejati --- kita memakainya terus-menerus
di bawah ini.
\end{example}

\section{Keterbilangan dan kardinalitas}

\begin{definition}[Ekuipotensi, keterbilangan]\label{def:b2:structures:countable}
Dua himpunan disebut \emph{ekuipoten}\index{ekuipotensi} apabila ada
bijeksi yang menghubungkannya. Suatu himpunan disebut
\emph{terbilang}\index{himpunan terbilang} apabila ekuipoten dengan $\N$
(sebagian penulis memasukkan himpunan hingga; kita memakai ungkapan
\emph{paling banyak terbilang} untuk ``hingga atau terbilang'').
\end{definition}

\begin{proposition}[Sifat kestabilan]\label{prop:b2:structures:countablestable}
\begin{enumerate}
  \item Setiap himpunan bagian tak hingga dari $\N$ adalah \omterm{def:b2:structures:countable}{terbilang};
        suatu himpunan paling banyak \omterm{def:b2:structures:countable}{terbilang} bila dan hanya bila ia
        terinjeksi ke $\N$, bila dan hanya bila ia kosong atau
        merupakan peta surjektif dari $\N$.
  \item $\N \times \N$ \omterm{def:b2:structures:countable}{terbilang}; hasil kali dua himpunan yang paling
        banyak \omterm{def:b2:structures:countable}{terbilang} juga paling banyak \omterm{def:b2:structures:countable}{terbilang}.
  \item Gabungan yang paling banyak \omterm{def:b2:structures:countable}{terbilang} atas himpunan yang paling
        banyak \omterm{def:b2:structures:countable}{terbilang} tetap paling banyak \omterm{def:b2:structures:countable}{terbilang}.
  \item $\Z$ dan $\Q$ \omterm{def:b2:structures:countable}{terbilang}.
\end{enumerate}
\end{proposition}

\begin{proof}
(1) Susun $A \subseteq \N$ yang tak hingga lewat minimum berulang: $a_0 =
\min A$, $a_{k+1} = \min\,(A \setminus \{a_0, \dots, a_k\})$
(tak kosong karena $A$ tak hingga); pemetaan $k \mapsto a_k$ naik tegas,
injektif, dan surjektif ke $A$ (setiap $a \in A$ hanya melampaui
berhingga banyak unsur $A$, jadi pasti tercapai). Jika $E$ terinjeksi ke
$\N$ lewat $\varphi$, maka $E$ \omterm{def:b2:structures:countable}{ekuipoten} dengan $\varphi(E) \subseteq
\N$: hingga atau \omterm{def:b2:structures:countable}{terbilang}. Jika $s \colon \N \to E$ surjektif, maka $x
\mapsto \min s^{-1}(\{x\})$ menginjeksikan $E$ ke $\N$.

(2) Pemetaan $(p, q) \mapsto 2^p(2q + 1) - 1$ adalah bijeksi $\N^2
\to \N$ (setiap bilangan bulat positif punya pemisahan ganjil--genap
tunggal $2^p m$ dengan $m$ ganjil, menurut ketunggalan pemfaktoran).
Untuk hasil kali: susun kedua injeksi.

(3) Diberikan himpunan $E_n$ beserta surjeksi $s_n \colon \N \to E_n$
(tak masalah bila ada $E_n$ yang hingga: ulangi saja nilainya), pemetaan
$(n, k) \mapsto s_n(k)$ adalah surjeksi dari $\N^2$ yang \omterm{def:b2:structures:countable}{terbilang} ke
$\bigcup E_n$.

(4) $\Z = \N \cup (-\N^*)$: gabungan \omterm{def:b2:structures:countable}{terbilang}. $\Q$ adalah peta
surjektif dari $\Z \times \N^*$ (pemetaan pecahan), jadi paling banyak
\omterm{def:b2:structures:countable}{terbilang}, sekaligus tak hingga.
\end{proof}

\begin{example}[Fungsi pemasangan, dihitung]\label{ex:b2:structures:pairing}
Bijeksi $(p, q) \mapsto 2^p(2q + 1) - 1$ dalam bukti di atas
pantas dilihat sedang bekerja. Nilai-nilai pertamanya:
\[
\begin{array}{c|ccccc}
 & q = 0 & q = 1 & q = 2 & q = 3 & q = 4\\
\hline
p = 0 & 0 & 2 & 4 & 6 & 8\\
p = 1 & 1 & 5 & 9 & 13 & 17\\
p = 2 & 3 & 11 & 19 & 27 & 35\\
p = 3 & 7 & 23 & 39 & 55 & 71
\end{array}
\]
Baris $p$ menghimpun bilangan $n$ yang $n + 1$-nya habis dibagi
tepat oleh $2^p$: setiap bilangan cacah muncul tepat sekali.
Pembacaan baliknya segamblang penyandiannya: untuk $n = 43$,
faktorkan $n + 1 = 44 = 2^2\cdot 11 = 2^2(2\cdot5 + 1)$, sehingga
$(p, q) = (2, 5)$. Pelajaran penutupnya: bukti keterbilangan
sering kali berupa \emph{algoritme} yang menyamar --- di sini,
``keluarkan semua faktor dua''.
\end{example}

\begin{example}[Bilangan aljabar adalah terbilang]\label{ex:b2:structures:algebraic}
Suatu bilangan kompleks disebut \emph{\omterm{def:b2:structures:algebra}{aljabar}} apabila ia menolkan
suatu polinomial tak nol berkoefisien rasional. Himpunan
$\overline\Q$ semua bilangan \omterm{def:b2:structures:algebra}{aljabar} adalah \omterm{def:b2:structures:countable}{terbilang}: polinomial
berderajat $\leq d$ atas $\Q$ terinjeksi ke $\Q^{d+1}$, yaitu hasil
kali hingga atas \omterm{def:b2:structures:countable}{himpunan terbilang}
(\cref{prop:b2:structures:countablestable}~(2)); gabungan atas
$d$ menyusun semua polinomial rasional tak nol sebagai $P_0, P_1,
P_2, \dots$; setiap $P_k$ punya berhingga banyak akar; dan
\[
\overline\Q = \bigcup_{k \in \N}\ \{\text{akar-akar } P_k\}
\]
adalah gabungan \omterm{def:b2:structures:countable}{terbilang} atas himpunan hingga
(\cref{prop:b2:structures:countablestable}~(3)), tak hingga karena
memuat $\Q$. Digabungkan dengan ketakterbilangan $\R$
(\cref{thm:b2:structures:cantor} di bawah), ini membuktikan ---
tanpa menunjukkan satu pun contohnya --- bahwa bilangan transenden
memang ada dan justru merupakan mayoritas yang tak \omterm{def:b2:structures:countable}{terbilang}:
itulah hujah pencacahan Cantor tahun 1874, keberadaan semata-mata
lewat kardinalitas.
\end{example}

\begin{theorem}[Cantor; ketakterbilangan $\R$]\label{thm:b2:structures:cantor}
\begin{enumerate}
  \item Untuk setiap himpunan $E$, tidak ada surjeksi $E \to
        \mathcal{P}(E)$.
  \item $\R$ \emph{tidak} \omterm{def:b2:structures:countable}{terbilang}.
\end{enumerate}
\end{theorem}

\begin{proof}
(1) telah dibuktikan pada jilid Tahun ke-1 (lewat himpunan diagonal $D =
\{x : x \notin f(x)\}$).

(2) Andaikan $(x_n)_{n \in \N}$ mendaftar seluruh $\R$. Bangun ruas
bersarang $I_0 \supseteq I_1 \supseteq \dots$ dengan $\abs{I_n} = 3^{-n}$
dan $x_n \notin I_n$: bagi ruas yang sedang dipegang menjadi tiga
pertigaan tertutup; paling sedikit satu pertigaan menghindari $x_n$
(sebuah titik menyinggung paling banyak dua dari ketiganya). Teorema
ruas bersarang (ujung-ujung yang berdampingan) memberi $\ell \in
\bigcap_n I_n$; tetapi $\ell = x_N$ untuk suatu $N$, padahal $x_N
\notin I_N$: kontradiksi.
\end{proof}

\begin{theorem}[Cantor--Bernstein]\label{thm:b2:structures:cantorbernstein}
Jika $E$ terinjeksi ke $F$ dan $F$ terinjeksi ke $E$, maka $E$ dan $F$
\omterm{def:b2:structures:countable}{ekuipoten}.\index{teorema Cantor--Bernstein}
\end{theorem}

\begin{proof}
Misalkan $f \colon E \to F$ dan $g \colon F \to E$ dua injeksi. Untuk
setiap titik (di $E$ maupun di $F$), telusuri \emph{rantai leluhurnya},
yakni barisan prapeta berurutan $x \mapsto g^{-1}(x) \mapsto
f^{-1}(g^{-1}(x)) \mapsto \dots$ --- setiap langkah terdefinisi selama
titik yang sedang dipegang terletak di peta injeksi yang bersangkutan,
dan langkah itu tunggal berkat keinjektifan. Ada tiga nasib yang saling
lepas: rantai berhenti di suatu titik $E \setminus g(F)$ (\emph{berpangkal
di $E$}), berhenti di suatu titik $F \setminus f(E)$ (\emph{berpangkal di
$F$}), atau tak pernah berhenti. Ini memilah $E = E_E \cup E_F \cup
E_\infty$ dan $F = F_E \cup F_F \cup F_\infty$ menurut pangkalnya.

Sekarang amati: $f$ memetakan $E_E$ \emph{pada} $F_E$ --- rantai
$f(x)$ adalah rantai $x$ yang didahului satu langkah, jadi pangkalnya
sama; dan setiap $y \in F_E$ punya rantai dengan sekurang-kurangnya satu
langkah (pangkalnya di $E$), sehingga $y = f(x)$ dengan $x \in E_E$.
Alasan yang sama memberi bijeksi $f \colon E_\infty \to F_\infty$ dan $g
\colon F_F \to E_F$. Setelah direkatkan,
\[
h(x) =
\begin{cases}
f(x) & \text{jika } x \in E_E \cup E_\infty,\\
g^{-1}(x) & \text{jika } x \in E_F,
\end{cases}
\]
merupakan bijeksi dari $E$ pada $F = F_E \cup F_\infty \cup F_F$: ia
bijektif sepotong demi sepotong, dan ketiga potongan sasarannya saling lepas.
\end{proof}

\begin{example}\label{ex:b2:structures:cbexample}
$\intoo{0}{1}$ dan $\intcc{0}{1}$ \omterm{def:b2:structures:countable}{ekuipoten}: identitas menginjeksikan
satu arah, $x \mapsto \frac{x + 1}{3}$ arah sebaliknya; teorema tadi
memproduksi bijeksinya (yang mau tak mau tak kontinu). Demikian pula
$\R$, $\intoo{0}{1}$ (lewat bijeksi bertipe $\tanh$) dan
$\mathcal{P}(\N)$ (uraian biner, \cref{exo:b2:structures:3})
semuanya \omterm{def:b2:structures:countable}{ekuipoten}: itulah ``kardinalitas kontinum''.
\end{example}

\begin{example}[Ruas dan persegi]\label{ex:b2:structures:squareseg}
$\intcc{0}{1}$ dan $\intcc{0}{1}^2$ \omterm{def:b2:structures:countable}{ekuipoten} --- dimensi tidak
tertangkap oleh kardinalitas. Satu injeksinya sepele: $x
\mapsto (x, 0)$. Untuk arah sebaliknya, kirim $(x, y)$ ke bilangan real
yang angka desimalnya menyelang-nyeling angka $x$ dan angka $y$,
\[
(0.x_1x_2x_3\dots,\ 0.y_1y_2y_3\dots)
\;\longmapsto\; 0.x_1y_1x_2y_2x_3y_3\dots,
\]
dengan memilih untuk tiap koordinat uraian yang tidak berakhir
dengan $9$ berulang: dengan kesepakatan itu angka pada petanya
menentukan angka $x$ dan $y$, jadi pemetaannya injektif (ia tidak
harus surjektif --- peta tidak pernah memuat, misalnya, bilangan
yang angka-angka pada posisi ganjilnya akhirnya selalu $9$ --- dan
itu tidak menjadi soal). Cantor--Bernstein
(\cref{thm:b2:structures:cantorbernstein}) lalu merakit bijeksi
sejati. Kekontinuan, tentu saja, mustahil: tidak ada bijeksi kontinu
di antara keduanya --- bab-bab metrik menjelaskan sebabnya
(keterhubungan membedakan garis dari bidang, \cref{ch:b2:metric}).
\end{example}

\section{Grup}

\begin{definition}[Subgrup yang dibangun; orde]\label{def:b2:structures:generated}
Misalkan $G$ grup dan $A \subseteq G$. Subgrup yang \emph{dibangun}
oleh $A$, ditulis $\langle A \rangle$, adalah subgrup terkecil yang
memuat $A$ --- secara konkret, semua hasil kali hingga atas unsur $A$
dan inversnya. Suatu grup disebut \emph{siklik}\index{grup siklik}
apabila dibangun oleh satu unsur: $\langle a\rangle = \{a^k : k \in \Z\}$.
Adapun \emph{orde}\index{orde!suatu unsur} unsur $a \in G$ adalah
$\operatorname{ord}(a) = \abs{\langle a \rangle}$ (boleh jadi tak
hingga); bila hingga, ia adalah $n \geq 1$ terkecil dengan $a^n = e$,
dan $a^k = e \iff \operatorname{ord}(a) \mid k$.
\end{definition}

\begin{proof}[Bukti pencirian orde]
Jika ada $a^m = e$ dengan $m \geq 1$, ambil $n \geq 1$ terkecil dengan
$a^n = e$. Unsur $e, a, \dots, a^{n-1}$ berbeda sepasang demi sepasang
($a^{i} = a^{j}$ dengan $0 \leq i < j < n$ memberi $a^{j-i} = e$, yang
melawan keminimalan), dan setiap $a^k$ menyusut ke salah satunya lewat
pembagian Euklides $k = nq + r$: jadi $\langle a\rangle$ tepat punya $n$
unsur, dan $a^k = a^r = e \iff r = 0 \iff n \mid k$. Jika tidak ada
pangkat yang trivial, semua $a^k$ ($k \in \Z$) berbeda (dengan alasan
pembagian yang sama) dan \omterm{def:b2:structures:generated}{ordenya} tak hingga.
\end{proof}

\begin{theorem}[Lagrange]\label{thm:b2:structures:lagrange}
Misalkan $G$ grup hingga dan $H$ subgrupnya. Maka $\abs H$ membagi
$\abs G$.\index{teorema Lagrange} Khususnya \omterm{def:b2:structures:generated}{orde} setiap unsur membagi
$\abs G$, dan $a^{\abs G} = e$ untuk setiap $a \in G$.
\end{theorem}

\begin{proof}
Relasi $x \sim y \iff x^{-1}y \in H$ adalah relasi ekuivalensi
(refleksif: $e \in H$; simetris: lewat invers; transitif: lewat hasil
kali). Kelas $x$ adalah \emph{koset kiri} $xH = \{xh : h \in H\}$,
dan $h \mapsto xh$ merupakan bijeksi $H \to xH$ (dengan invers $y \mapsto
x^{-1}y$): jadi semua kelas punya $\abs H$ unsur. Kelas-kelas itu
memilah $G$ (teorema pemilahan umum pada jilid Tahun ke-1), sehingga
$\abs G = \abs H \times (\text{banyaknya koset})$. Untuk sebuah unsur:
terapkan hasil ini pada $H = \langle a\rangle$; maka $a^{\abs G} =
(a^{\operatorname{ord} a})^{\abs G / \operatorname{ord} a} = e$.
\end{proof}

\begin{example}[Koset dalam kerja: $A_3$ di dalam $\mathfrak{S}_3$]\label{ex:b2:structures:cosets}
Ambil $G = \mathfrak{S}_3$ (berorde $6$) dan $H = A_3 =
\{\mathrm{id},\ (1\,2\,3),\ (1\,3\,2)\}$. Koset kirinya adalah
\[
H = \{\mathrm{id},\ (1\,2\,3),\ (1\,3\,2)\},
\qquad
(1\,2)H = \{(1\,2),\ (2\,3),\ (1\,3)\} :
\]
dua kelas beranggotakan tiga unsur yang memilah $G$, persis seperti
dituntut oleh pencacahan $\abs G = \abs H \times (\text{banyaknya koset})$
--- dan tampak jelas bahwa itulah pemilahan atas permutasi genap dan
permutasi ganjil. Perhatikan $(1\,3)H = (1\,2)H$ walaupun $(1\,3) \neq
(1\,2)$: koset adalah \emph{kelas}, bukan sesuatu yang dinamai menurut
wakilnya, dan $x^{-1}y \in H$ satu-satunya perbandingan yang sah.
Gambaran dua kelas ini berlaku umum untuk tanda permutasi: $A_n$
bersama satu-satunya koset pendampingnya membelah $\mathfrak{S}_n$ tepat
menjadi dua, dan begitulah soal akhir pekan bab ini mencacah posisi
teka-teki yang terjangkau.
\end{example}

\begin{example}\label{ex:b2:structures:lagrangeapp}
Dua panen langsung. \emph{Grup berorde prima pasti \omterm{def:b2:structures:generated}{siklik}:}
jika $\abs G = p$ prima dan $a \neq e$, maka
$\operatorname{ord}(a)$ membagi $p$ dan tidak sama dengan $1$, jadi
ia sama dengan $p$: $\langle a\rangle = G$. \emph{Kisi subgrup
$\Z/12\Z$:} menurut \cref{prop:b2:structures:cyclic} di bawah, ada
tepat satu subgrup untuk tiap pembagi $12$ --- berorde $1, 2, 3, 4,
6, 12$, yang berturut-turut \omterm{def:b2:structures:generated}{dibangun} oleh $\overline 0$, $\overline
6$, $\overline 4$, $\overline 3$, $\overline 2$, $\overline 1$.
Peringatan penutup: \emph{konvers} teorema Lagrange tidak berlaku
umum --- $A_4$ berorde $12$ tetapi tidak punya subgrup berorde $6$,
sebagaimana kita buktikan pada soal akhir pekan bab ini
(\cref{pb:b2:structures:1}, pertanyaan 14). Lagrange membatasi \omterm{def:b2:structures:generated}{orde}
yang mungkin; ia tidak menjanjikan \omterm{def:b2:structures:generated}{orde} itu ada.
\end{example}

\begin{omfigure}
\begin{tikzpicture}[scale=1.05, every node/.style={font=\small}]
  \node (G) at (0,3)  {$\Z/12\Z = \langle\overline1\rangle$};
  \node (A) at (-1.6,2) {$\langle\overline2\rangle$};
  \node (B) at (1.6,2)  {$\langle\overline3\rangle$};
  \node (C) at (-1.6,1) {$\langle\overline4\rangle$};
  \node (D) at (1.6,1)  {$\langle\overline6\rangle$};
  \node (E) at (0,0)  {$\{\overline0\}$};
  \draw[omDef] (G) -- (A); \draw[omDef] (G) -- (B);
  \draw[omDef] (A) -- (C); \draw[omDef] (A) -- (D);
  \draw[omDef] (B) -- (D); \draw[omDef] (C) -- (E);
  \draw[omDef] (D) -- (E);
  \node[omThm, right=2pt] at (0.1,3) {\ orde $12$};
  \node[omThm, left=2pt]  at (-2.1,2) {orde $6$\ };
  \node[omThm, right=2pt] at (2.1,2) {\ orde $4$};
  \node[omThm, left=2pt]  at (-2.1,1) {orde $3$\ };
  \node[omThm, right=2pt] at (2.1,1) {\ orde $2$};
\end{tikzpicture}

{\small Kisi subgrup $\Z/12\Z$: satu subgrup untuk tiap
pembagi $12$ (\cref{prop:b2:structures:cyclic}), dengan sebuah
rusuk bila yang satu memuat yang lain dengan indeks prima.
Pemuatan berjalan \emph{berlawanan} dengan keterbagian pembangunnya:
$\langle\overline 4\rangle \subseteq \langle\overline2\rangle$ karena
$4$ adalah kelipatan $2$.}
\end{omfigure}

\begin{proposition}[Grup siklik]\label{prop:b2:structures:cyclic}
Misalkan $G = \langle a \rangle$ \omterm{def:b2:structures:generated}{siklik} berorde $n$.
\begin{enumerate}
  \item $G$ isomorfik dengan $(\Z/n\Z, +)$, lewat $\overline k \mapsto
        a^k$.
  \item Setiap subgrup $G$ adalah \omterm{def:b2:structures:generated}{siklik}; untuk tiap pembagi $d \mid n$
        ada tepat satu subgrup berorde $d$, yakni $\langle
        a^{n/d}\rangle$.
  \item $a^k$ membangun $G$ bila dan hanya bila $\gcd(k, n) = 1$: jadi
        $G$ punya $\varphi(n)$ pembangun (fungsi Euler\index{fungsi
        Euler}).
\end{enumerate}
\end{proposition}

\begin{proof}
(1) Pemetaan $k \mapsto a^k$ dari $\Z$ pada $G$ selaras dengan
kekongruenan modulo $n$ ($a^{k} = a^{k'} \iff n \mid k - k'$, menurut
pencirian \omterm{def:b2:structures:generated}{orde}); sifat universal
(\cref{def:b2:structures:quotient}) menghasilkan morfisma bijektif yang
terdefinisi dengan baik dari $\Z/n\Z$.

(2) Misalkan $H \leq G$ tak trivial dan $m$ bilangan $\geq 1$ terkecil
dengan $a^m \in H$. Pembagian Euklides menunjukkan $H = \langle
a^m\rangle$ (untuk $a^k \in H$: dari $k = mq + r$ terpaksa $a^r \in H$,
jadi $r = 0$), dan $m \mid n$ (bagilah $n$ oleh $m$: $a^{n \bmod m} \in
H$). Akibatnya $\abs H = n/m$; dengan mengambil $m = n/d$ setiap pembagi
$d$ terwujud. Ketunggalan: menurut uraian di atas, sembarang subgrup
berorde $d$ berbentuk $\langle a^m \rangle$ dengan $n/m = d$ ---
sehingga $m = n/d$ terpaksa dan subgrupnya tertentu.

(3) Kita klaim $\operatorname{ord}(a^k) = \frac{n}{\gcd(k, n)}$.
Tulis $d = \gcd(k, n)$. Untuk sembarang $m \geq 1$, pencirian \omterm{def:b2:structures:generated}{orde} pada
\cref{def:b2:structures:generated} memberi rantai kesetaraan berikut
\[
(a^k)^m = e
\iff n \mid km
\iff \frac{n}{d} \,\Big|\, \frac{k}{d}\,m
\iff \frac{n}{d} \,\Big|\, m ,
\]
langkah terakhir memakai lema Gauss, karena $\frac nd$ dan $\frac kd$
saling prima. Bilangan $m$ terkecil yang demikian adalah $\frac nd$:
$\operatorname{ord}(a^k) = \frac n{\gcd(k,n)}$, yang sama dengan $n$
bila dan hanya bila $\gcd(k, n) = 1$. Ada $\varphi(n)$ kelas $k$
modulo $n$ yang seperti itu.
\end{proof}

\section{Grup simetri}

\begin{definition}\label{def:b2:structures:sn}
$\mathfrak{S}_n$ adalah grup permutasi atas $\intint{1}{n}$
(berorde $n!$). Sebuah \emph{siklus}\index{siklus} $(a_1\,a_2\,\cdots\,a_k)$
memetakan $a_1 \mapsto a_2 \mapsto \dots \mapsto a_k \mapsto a_1$ dan
membiarkan yang lain tetap; $k$ disebut \emph{panjangnya}, dan siklus
berpanjang $2$ disebut \emph{transposisi}\index{transposisi}. Dua siklus
disebut \emph{saling lepas} apabila penyangganya (titik yang tidak
tetap) saling lepas.
\end{definition}

\begin{theorem}[Penguraian siklus]\label{thm:b2:structures:cycles}
Setiap permutasi $\sigma \neq \mathrm{id}$ adalah hasil kali \omterm{def:b2:structures:sn}{siklus}
yang saling lepas, tunggal kecuali urutan faktornya. \omterm{def:b2:structures:sn}{Siklus} yang saling
lepas komutatif, dan $\operatorname{ord}(\sigma)$ adalah KPK
panjang-panjangnya.
\end{theorem}

\begin{proof}
Tinjau relasi ``orbit'' pada penyangga $\sigma$: $x \sim y$ bila
$y = \sigma^k(x)$ untuk suatu $k \in \Z$ --- sebuah relasi
ekuivalensi. Setiap kelasnya berbentuk
$\{x, \sigma(x), \dots, \sigma^{k-1}(x)\}$
(hingga, jadi iterasinya pasti berputar kembali --- pengulangan pertama
mesti kembali ke $x$ berkat keinjektifan) mengusung \omterm{def:b2:structures:sn}{siklus} $(x\
\sigma(x)\ \cdots\ \sigma^{k-1}(x))$, dan $\sigma$ adalah hasil kali
\omterm{def:b2:structures:sn}{siklus-siklus} itu: pada tiap orbit hanya \omterm{def:b2:structures:sn}{siklus} yang bersangkutan yang
bekerja. Ketunggalan: sembarang pemfaktoran atas \omterm{def:b2:structures:sn}{siklus} saling lepas
melahirkan orbit yang persis sama (\omterm{def:b2:structures:sn}{siklus} lewat $x$ mestilah $(x\
\sigma(x)\ \cdots)$). \omterm{def:b2:structures:sn}{Siklus} saling lepas komutatif karena menggerakkan
titik yang berbeda; pernyataan tentang \omterm{def:b2:structures:generated}{orde} menyusul karena $\sigma^m =
\mathrm{id}$ bila dan hanya bila pangkat ke-$m$ tiap \omterm{def:b2:structures:sn}{siklus} demikian
(karena saling lepas), bila dan hanya bila tiap panjang membagi $m$.
\end{proof}

\begin{example}[Tipe siklus sebagai pencacahan]\label{ex:b2:structures:cycletype}
Ada berapa permutasi di $\mathfrak{S}_9$ bertipe \omterm{def:b2:structures:sn}{siklus}
$(4, 3, 2)$ --- satu \omterm{def:b2:structures:sn}{siklus} panjang $4$, satu \omterm{def:b2:structures:sn}{siklus} panjang $3$, dan
satu \omterm{def:b2:structures:sn}{transposisi}? Pilih penyangganya sekaligus urutan siklisnya:
\[
\frac{9!}{4\cdot 3\cdot 2}
= \frac{362\,880}{24} = 15\,120 :
\]
jajarkan kesembilan lambang itu dalam satu baris ($9!$ cara), kurung
empat yang pertama, tiga berikutnya, dan dua terakhir menjadi
\omterm{def:b2:structures:sn}{siklus}, lalu bagi dengan banyaknya perputaran di dalam tiap kurung
($4$, $3$ dan $2$ buah) yang memberi permutasi yang sama.
(\emph{Panjang} siklusnya di sini berbeda-beda, jadi tidak perlu
pembagian lagi; panjang yang sama akan menuntut pembagian oleh
permutasi antar kurung yang sepanjang itu pula.) Setiap permutasi
semacam itu berorde $\operatorname{lcm}(4,3,2) = 12$ dan bertanda
$(-1)^3(-1)^2(-1)^1 = +1$
(\cref{thm:b2:structures:cycles} dan teorema tanda permutasi di
bawah). Satu partisi $9$, satu kelas konjugasi, satu pencacahan
--- kombinatorika $\mathfrak{S}_n$ tak lain adalah aritmetika
partisi.
\end{example}

\begin{theorem}[Tanda permutasi]\label{thm:b2:structures:signature}
Hanya ada satu morfisma grup $\varepsilon \colon
\mathfrak{S}_n \to \{\pm 1\}$ (untuk $n \geq 2$) yang bernilai $-1$
pada \omterm{def:b2:structures:sn}{transposisi}: itulah \emph{tanda permutasi}\index{tanda permutasi}.
Lebih lanjut $\varepsilon(\sigma) = (-1)^{I(\sigma)}$ dengan $I(\sigma)$
menyatakan banyaknya \emph{inversi} (pasangan $i < j$ dengan $\sigma(i) >
\sigma(j)$), \omterm{def:b2:structures:sn}{siklus} berpanjang $k$ bertanda $(-1)^{k-1}$, dan
\emph{grup alternating} $A_n = \ker\varepsilon$ berorde
$\frac{n!}{2}$.
\end{theorem}

\begin{proof}
\emph{Keberadaan.} Untuk $\sigma \in \mathfrak{S}_n$ tetapkan
\[
\varepsilon(\sigma)
= \prod_{1 \leq i < j \leq n}
\frac{\sigma(j) - \sigma(i)}{j - i} .
\]
Nilai mutlak faktor-faktornya berkali menjadi $1$ (pasangan tak terurut
$\{\sigma(i), \sigma(j)\}$ menjelajahi semua pasangan), sehingga
$\varepsilon(\sigma) = (-1)^{I(\sigma)} \in \{\pm1\}$. Morfisma: untuk
$\sigma, \tau$,
\[
\varepsilon(\sigma\tau)
= \prod_{i<j} \frac{\sigma(\tau(j)) - \sigma(\tau(i))}{j - i}
= \prod_{i<j} \frac{\sigma(\tau(j)) - \sigma(\tau(i))}{\tau(j) -
\tau(i)} \cdot \prod_{i<j} \frac{\tau(j) - \tau(i)}{j - i}
= \varepsilon(\sigma)\,\varepsilon(\tau),
\]
karena hasil kali di tengah sama dengan $\varepsilon(\sigma)$ setelah
diindeks ulang menurut pasangan $\{\tau(i), \tau(j)\}$ (tiap pasangan
tak terurut muncul sekali, dan pembilang serta penyebut berganti tanda
bersama-sama). \omterm{def:b2:structures:sn}{Transposisi} $\tau = (a\,b)$ dengan $a < b$ punya
inversi sebanyak bilangan ganjil; dicacah dengan saksama: pasangan
terbalik $(i, j)$, $i < j$, dengan $\tau(i) > \tau(j)$ adalah
\[
(a, j) \ \text{untuk } a < j < b, \qquad
(i, b) \ \text{untuk } a < i < b, \qquad
(a, b) \ \text{sendiri},
\]
yakni sebanyak $(b - a - 1) + (b - a - 1) + 1 = 2(b - a) - 1$,
sebuah bilangan ganjil. (Cara lain: periksa langsung $(1\,2)$, yang
punya satu inversi, lalu konjugasikan --- unsur yang sekonjugasi
bertanda sama karena $\varepsilon$ adalah morfisma ke grup abelian.)
Jadi $\varepsilon((a\,b)) = (-1)^{2(b-a)-1} = -1$.

\emph{Ketunggalan.} \omterm{def:b2:structures:sn}{Transposisi} membangun $\mathfrak{S}_n$ (sebab untuk
sebarang \omterm{def:b2:structures:sn}{siklus} berlaku
$(a_1\cdots a_k) = (a_1\,a_k)(a_1\,a_{k-1})\cdots(a_1\,a_2)$,
lalu \cref{thm:b2:structures:cycles} merampungkannya); morfisma ke
$\{\pm1\}$ ditentukan oleh nilainya pada pembangun.

\emph{Akibat.} Kesamaan \omterm{def:b2:structures:sn}{siklus} di atas menuliskan \omterm{def:b2:structures:sn}{siklus} berpanjang $k$
sebagai $k - 1$ \omterm{def:b2:structures:sn}{transposisi}: tandanya $(-1)^{k-1}$. Tentang $A_n$:
morfisma $\varepsilon$ surjektif (\omterm{def:b2:structures:sn}{transposisi} ada untuk $n \geq 2$),
dan kedua ``koset'' $A_n$ dan $(1\,2)A_n$ \omterm{def:b2:structures:countable}{ekuipoten} serta memilah
$\mathfrak{S}_n$ (dengan hujah Lagrange): jadi $\abs{A_n} =
\frac{n!}{2}$.
\end{proof}

\begin{example}\label{ex:b2:structures:snexample}
$\sigma = \begin{pmatrix} 1&2&3&4&5&6\\ 3&6&5&4&1&2 \end{pmatrix}
= (1\,3\,5)(2\,6)$: berorde $\operatorname{lcm}(3,2) = 6$, bertanda
$(-1)^{2}\cdot(-1)^{1} = -1$. Tanda permutasi adalah pemeriksaan
paritas tercepat atas suatu pengocokan --- sekaligus mesin penggerak
determinan pada \cref{ch:b2:linalg}.
\end{example}

\begin{example}[Tiga jalan menuju satu tanda]\label{ex:b2:structures:threeroads}
Misalkan $\sigma \in \mathfrak{S}_5$ mengirim $1, 2, 3, 4, 5$ ke $3, 5, 4,
1, 2$. \emph{Lewat \omterm{def:b2:structures:sn}{siklus}:} $1 \mapsto 3 \mapsto 4 \mapsto 1$ dan $2
\mapsto 5 \mapsto 2$, jadi $\sigma = (1\,3\,4)(2\,5)$ dan
$\varepsilon(\sigma) = (-1)^{2}(-1)^{1} = -1$. \emph{Lewat
inversi:} pada deretan nilai $3, 5, 4, 1, 2$ pasangan yang terbalik
adalah $(3,1)$, $(3,2)$, $(5,4)$, $(5,1)$, $(5,2)$, $(4,1)$,
$(4,2)$: tujuh buah, dan $(-1)^7 = -1$. \emph{Lewat
\omterm{def:b2:structures:sn}{transposisi}:} $\sigma = (1\,4)(1\,3)(2\,5)$, tiga faktor,
$(-1)^3 = -1$. Tiga perhitungan, satu paritas: ketunggalan pada
\cref{thm:b2:structures:signature} menjamin tidak ada tata cara
pembukuan yang dapat membuat ketiganya berselisih --- dan justru itulah
yang membuat $\varepsilon$ berguna sebagai invarian (lihat soal akhir pekan).
\end{example}

\begin{remark}[Ke mana tanda permutasi melangkah setelah ini]
Tanda permutasi adalah benih tiga panen berikutnya: ia membangun
determinan beserta aturan hasil kalinya pada \cref{ch:b2:linalg}; ia
menggerakkan invarian paritas untuk teka-teki kombinatorial (soal akhir
pekan bab ini menyelesaikan teka-teki lima belas dengannya); dan grup
alternating $A_n$ yang didefinisikannya menjadi tokoh utama pada jilid
Tahun ke-3, tempat kesederhanaannya untuk $n \geq 5$ menjelaskan
mengapa persamaan berderajat $5$ tidak terpecahkan dengan akar.
\end{remark}

\section{Ring, ideal, kuosien}

\begin{definition}[Ideal]\label{def:b2:structures:ideal}
Misalkan $A$ ring komutatif. Sebuah \emph{ideal}\index{ideal} $I
\subseteq A$ adalah subgrup aditif dengan sifat $a x \in I$ untuk setiap
$a \in A$, $x \in I$. Kernel morfisma ring adalah ideal; $I = A$ bila dan
hanya bila $1 \in I$, bila dan hanya bila $I$ memuat sebuah unit. Ideal
yang \emph{\omterm{def:b2:structures:generated}{dibangun}} oleh $x$ adalah $xA = \{xa\}$ (sebuah ideal
\emph{utama}).
\end{definition}

\begin{theorem}[{Ideal pada $\Z$ dan pada $K[X]$}]\label{thm:b2:structures:principal}
Setiap \omterm{def:b2:structures:ideal}{ideal} $\Z$ berbentuk $n\Z$ untuk suatu $n \in \N$ yang tunggal;
setiap \omterm{def:b2:structures:ideal}{ideal} $K[X]$ ($K$ sebuah lapangan) berbentuk $P\,K[X]$ untuk suatu
$P$ monik (atau nol) yang tunggal. Akibatnya FPB ada pada kedua ring itu
beserta relasi Bézout: $x\Z + y\Z = \gcd(x,y)\Z$, demikian pula untuk
polinomial.
\end{theorem}

\begin{proof}
Untuk $\Z$ ini tak lain teorema subgrup pada jilid Tahun ke-1 (sebuah
\omterm{def:b2:structures:ideal}{ideal} khususnya adalah subgrup, dan $n\Z$ memang \omterm{def:b2:structures:ideal}{ideal}). Untuk $K[X]$:
misalkan $I \neq \{0\}$ sebuah \omterm{def:b2:structures:ideal}{ideal} dan $P \in I$ tak nol berderajat
minimal, dinormalkan menjadi monik. Untuk $F \in I$, pembagian Euklides
$F = PQ + R$ memberi $R = F - PQ \in I$ dengan $\deg R < \deg P$:
keminimalan memaksa $R = 0$, jadi $I = P\,K[X]$. Ketunggalan: dua
pembangun monik saling membagi. Pernyataan Bézout tak lain kesamaan
\omterm{def:b2:structures:ideal}{ideal} $x\Z + y\Z$ (atau padanan polinomialnya) dengan \omterm{def:b2:structures:ideal}{ideal} utama yang
\omterm{def:b2:structures:generated}{dibangun} FPB --- persis definisi FPB yang dipakai pada Tahun ke-1, kini
dikenali sebagai pernyataan tentang \omterm{def:b2:structures:ideal}{ideal}.
\end{proof}

\begin{example}[FPB polinomial, dua jalan]\label{ex:b2:structures:polygcd}
Hitung $\gcd(X^3 - 1,\ X^2 - 1)$ di $\Q[X]$. \emph{Lewat Euklides:}
\[
X^3 - 1 = X\,(X^2 - 1) + (X - 1),
\qquad
X^2 - 1 = (X + 1)(X - 1) + 0 ,
\]
jadi FPB-nya $X - 1$, dan penyulihan balik memberi relasi
Bézout
\[
X - 1 = 1\cdot(X^3 - 1) - X\cdot(X^2 - 1).
\]
\emph{Lewat \omterm{def:b2:structures:ideal}{ideal}:} \omterm{def:b2:structures:ideal}{ideal} $(X^3 - 1)\Q[X] + (X^2 - 1)\Q[X]$ bersifat
utama (\cref{thm:b2:structures:principal}); ia memuat $X - 1$
(lihat ungkapan di atas) dan termuat di $(X - 1)\Q[X]$ (kedua
pembangunnya nol di $1$, jadi keduanya kelipatan $X - 1$): maka
pembangun moniknya adalah $X - 1$. Pelajaran penutupnya: sudut
pandang \omterm{def:b2:structures:ideal}{ideal} mengenali FPB \emph{tanpa membagi} --- akar yang sama-sama
dimiliki keduanya menempatkan \omterm{def:b2:structures:ideal}{idealnya}, dan Euklides sekadar mengesahkannya.
\end{example}

\begin{definition}[Ring kuosien $\Z/n\Z$, ditinjau ulang]\label{def:b2:structures:quotientring}
Untuk sebuah \omterm{def:b2:structures:ideal}{ideal} $I$ pada $A$, relasi $x \sim y \iff x - y \in I$
adalah relasi ekuivalensi yang selaras dengan $+$ dan $\times$; \omterm{def:b2:structures:quotient}{himpunan
kuosien} $A/I$ mewarisi struktur ring --- itulah \emph{ring
kuosien}\index{ring kuosien} --- sehingga $\pi \colon A \to A/I$ menjadi
morfisma dengan kernel $I$. Untuk $A = \Z$, $I = n\Z$ ini tak lain
$\Z/n\Z$ pada jilid Tahun ke-1, kini lengkap dengan sifat universalnya:
setiap morfisma yang menolkan $I$ terfaktorkan lewat $A/I$.
\end{definition}

\begin{theorem}[Teorema sisa Cina, bentuk ring]\label{thm:b2:structures:crt}
Jika $\gcd(m, n) = 1$, maka pemetaan
\[
\Z/mn\Z \longrightarrow \Z/m\Z \times \Z/n\Z,
\qquad
\overline{x} \longmapsto (x \bmod m,\; x \bmod n)
\]
adalah isomorfisma ring.\index{teorema sisa Cina} Akibatnya
$\varphi(mn) = \varphi(m)\varphi(n)$ untuk $m, n$ saling prima, dan
\[
\varphi(n) = n \prod_{p \mid n} \Bigl(1 - \frac 1p\Bigr)
\quad (p \text{ prima}).
\]
\end{theorem}

\begin{proof}
Pemetaan itu morfisma ring yang terdefinisi dengan baik (keselarasannya
langsung terlihat). Keinjektifan: $x \equiv 0$ modulo $m$ dan modulo $n$
dengan $\gcd(m,n) = 1$ memaksa $mn \mid x$ (Gauss). Kesurjektifan: kedua
ruas punya $mn$ unsur, jadi keinjektifan sudah cukup (kardinalitas hingga
yang sama) --- atau secara gamblang: dari relasi Bézout $um +
vn = 1$, kelas
\[
x = b\,um + a\,vn
\]
terpetakan ke $(a \bmod m,\ b \bmod n)$, sebab $vn = 1 - um \equiv 1
\pmod m$ membuat $x \equiv a \pmod m$, dan setangkup untuk modulo $n$
--- itulah resep yang dipakai secara numerik pada
\cref{ex:b2:structures:crtinverse}. Unit berpadanan dengan pasangan unit
(unit sebuah ring hasil kali adalah pasangan unit), sehingga $\varphi(mn) =
\varphi(m)\varphi(n)$. Untuk pangkat prima, $\varphi(p^k) = p^k -
p^{k-1}$ (yang bukan unit modulo $p^k$ adalah kelipatan $p$);
kemultiplikatifan lalu merakit rumus hasil kalinya.
\end{proof}

\begin{example}[Membalik isomorfisma Cina]\label{ex:b2:structures:crtinverse}
Ambil $m = 8$, $n = 9$. Invers isomorfisma itu dibuat gamblang
lewat dua \emph{idempoten}: cari $u \equiv 1 \pmod 8$,
$u \equiv 0 \pmod 9$ serta $v \equiv 0 \pmod 8$, $v \equiv 1 \pmod
9$. Dari $u = 9k \equiv 1 \pmod 8$: $k \equiv 1$, jadi $u = 9$; dari
$v = 8k \equiv 1 \pmod 9$: $-k \equiv 1$, $k \equiv 8$, jadi $v =
64$. Maka kelas $x = 9a + 64b$ modulo $72$ adalah satu-satunya
penyelesaian $x \equiv a \pmod 8$, $x \equiv b \pmod 9$: untuk $a =
3$, $b = 5$ diperoleh $27 + 320 = 347 \equiv 59 \pmod{72}$ ---
persis nilai antara yang ditemukan lewat penyulihan pada
\cref{exo:b2:structures:8}. Pelajaran penutupnya: $u$ dan $v$
memenuhi $u + v \equiv 1$, $uv \equiv 0$, $u^2 \equiv u$, $v^2
\equiv v$ modulo $72$; keduanya adalah peta $(1, 0)$ dan $(0,
1)$, dan setiap penguraian Cina pada dasarnya adalah penguraian
$1$ menjadi idempoten yang saling ortogonal.
\end{example}

\begin{theorem}[Euler; Fermat ditinjau ulang]\label{thm:b2:structures:euler}
Unit-unit $\Z/n\Z$ membentuk grup berorde $\varphi(n)$; karenanya untuk
$\gcd(a, n) = 1$ berlaku
\[
a^{\varphi(n)} \equiv 1 \pmod n
\qquad (\text{teorema Euler\index{teorema Euler}}),
\]
dan teorema kecil Fermat adalah kasus $n = p$ prima, yang kini berjarak
satu baris dari teorema Lagrange.
\end{theorem}

\begin{proof}
Kelas yang punya invers persis kelas bilangan bulat yang saling prima
dengan $n$ (jilid Tahun ke-1): ada $\varphi(n)$ buah, dan semuanya
membentuk grup terhadap perkalian. Menurut Lagrange
(\cref{thm:b2:structures:lagrange}): setiap unsur dipangkatkan \omterm{def:b2:structures:generated}{orde}
grupnya menghasilkan unsur identitas.
\end{proof}

\begin{example}[Grup unit tanpa pembangun]\label{ex:b2:structures:units15}
Grup $(\Z/15\Z)^*$ punya $\varphi(15) = \varphi(3)\varphi(5)
= 8$ unsur. Apakah ia \omterm{def:b2:structures:generated}{siklik}? Hitung \omterm{def:b2:structures:generated}{ordenya} lewat isomorfisma
Cina $(\Z/15\Z)^* \simeq (\Z/3\Z)^* \times (\Z/5\Z)^*$
(unit modulo $15$ adalah pasangan unit): kedua faktornya berorde
$2$ dan $4$, jadi \omterm{def:b2:structures:generated}{orde} setiap unsur membagi
$\operatorname{lcm}(2, 4) = 4 < 8$ --- tidak ada unsur yang membangun.
Secara konkret:
\[
2^4 = 16 \equiv 1, \qquad
4^2 = 16 \equiv 1, \qquad
7^4 \equiv 1, \qquad
11^2 = 121 \equiv 1, \qquad
14^2 \equiv 1 \pmod{15} :
\]
\omterm{def:b2:structures:generated}{ordenya} $4, 2, 4, 2, 2$ dan tidak pernah $8$. Bandingkan dengan
\cref{exo:b2:structures:10}: $(\Z/p\Z)^*$ \emph{memang} \omterm{def:b2:structures:generated}{siklik} untuk
$p$ prima, sebab di sana grup unitnya berada di dalam sebuah lapangan.
Teorema Euler tetap berlaku dengan pangkat $\varphi(15) = 8$,
tetapi pangkat semesta yang sebenarnya di sini adalah $4$ --- Euler
memberi batas atas, tidak selalu batas yang tajam.
\end{example}

\begin{definition}[Aljabar]\label{def:b2:structures:algebra}
Sebuah \emph{aljabar-$K$}\index{aljabar} adalah ruang vektor $K$ bernama
$A$ yang dilengkapi struktur ring dengan perkalian yang bilinear atas
$K$. Contohnya: $K[X]$, $\mathcal{M}_n(K)$, $\mathcal{L}(E)$, ruang
fungsi $\mathcal{F}(X, K)$, dan $\C$ sebagai aljabar-$\R$. Morfisma
aljabar adalah morfisma ring yang linear; \emph{evaluasi} $P \mapsto P(u)$
dari $K[X]$ ke $\mathcal{L}(E)$ (atau ke $\mathcal{M}_n(K)$) adalah
contoh pusatnya, yang menggerakkan \cref{ch:b2:reduction}.
\end{definition}

\begin{example}[Morfisma evaluasi dan kernelnya]\label{ex:b2:structures:evaluation}
Ambil $A = \begin{pmatrix}0 & 1\\ 0 & 0\end{pmatrix}$ dan evaluasi
$\varepsilon_A \colon \R[X] \to \mathcal{M}_2(\R)$,
$P \mapsto P(A)$. Karena $A^2 = 0$,
\[
P(A) = P(0)\,I + P'(0)\,A =
\begin{pmatrix} P(0) & P'(0)\\ 0 & P(0)\end{pmatrix},
\]
(hanya suku tetap dan suku linear $P$ yang bertahan). Karenanya
$\ker\varepsilon_A = \{P : P(0) = P'(0) = 0\} = X^2\,\R[X]$: sebuah
\omterm{def:b2:structures:ideal}{ideal} utama, persis seperti diramalkan
\cref{thm:b2:structures:principal}, \omterm{def:b2:structures:generated}{dibangun} oleh $X^2$ yang monik
dan berderajat terkecil di dalam kernel --- itulah \emph{polinomial
minimal} $A$, bintang \cref{ch:b2:reduction}. Petanya adalah \omterm{def:b2:structures:algebra}{aljabar}
komutatif berdimensi dua $\{aI + bA\}$: morfisma evaluasi menciutkan
$\R[X]$ yang berdimensi tak hingga menjadi \omterm{def:b2:structures:algebra}{aljabar} kecil yang terhitungkan.
\end{example}

\begin{remark}[Pandangan ke depan: tiga melodi yang perlu disimak]\label{rem:b2:structures:perspectives}
Tiga gagasan struktural dari bab ini berulang sepanjang jilid ini,
setiap kali dengan orkestrasi yang makin tebal. \emph{Pemfaktoran
lewat kuosien} (\cref{def:b2:structures:quotient}): ia membangun
$\Z/n\Z$ di sini, mendefinisikan pemetaan pada himpunan penyelesaian
sistem linear di \cref{ch:b2:linalg}, dan diam-diam menopang setiap
hujah ``terdefinisi dengan baik pada kelas''. \emph{Invarian}: tanda
permutasi adalah morfisma ke $\{\pm1\}$ yang tak terelakkan oleh
langkah sah mana pun --- logika yang sama memberi aturan hasil kali
determinan (\cref{ch:b2:linalg}), keawetan trace terhadap keserupaan,
dan besaran kekal pada \cref{ch:b2:diffeq}.
\emph{Mencacah dengan bersandar pada struktur}: Lagrange mencacah lewat
koset, dimensi mencacah lewat basis
(\cref{ch:b2:linalg}), multiplisitas mencacah lewat derajat polinomial
(\cref{ch:b2:reduction}); setiap kali sebuah batas tampak ajaib, pasti
ada pemilahan atau penjenjangan yang sedang mencacah.
\end{remark}

\begin{remark}[Jebakan yang sering muncul]\label{rem:b2:structures:pitfalls}
Empat klasik. (i) Pemetaan pada suatu kuosien wajib diperiksa
\emph{keterdefinisiannya}: ``$\overline x \mapsto$ (rumus dalam $x$)''
sah hanya bila rumusnya tetap pada tiap kelas --- itulah keselarasan
pada \cref{def:b2:structures:quotient}, bukan sekadar
formalitas. (ii) $\operatorname{ord}(ab) =
\operatorname{lcm}(\operatorname{ord}a, \operatorname{ord}b)$
\emph{salah} secara umum, bahkan untuk unsur yang komutatif ($a$
dan $a^{-1}$); \cref{exo:b2:structures:4} memberi pernyataan yang
benar dengan syarat saling prima dan komutatif, sedangkan \omterm{def:b2:structures:sn}{siklus} saling
lepas memberi versi permutasinya. (iii) Keterbilangan bertahan
terhadap \emph{gabungan} \omterm{def:b2:structures:countable}{terbilang} dan \emph{hasil kali} hingga, tetapi
tidak terhadap hasil kali \omterm{def:b2:structures:countable}{terbilang}: $\{0,1\}^{\N}$ tidak \omterm{def:b2:structures:countable}{terbilang}
(\cref{exo:b2:structures:3}) walaupun tiap faktornya hanya punya dua
unsur. (iv) Cantor--Bernstein hanya menuntut injeksi dua arah,
tetapi bijeksi yang \omterm{def:b2:structures:generated}{dibangunnya} lazimnya tidak kontinu dan tidak
gamblang --- jangan berharap ada rumusnya
(\cref{ex:b2:structures:cbexample}).
\end{remark}

\begin{remark}[Di mana bab ini dipakai]
Hampir di mana-mana. Tanda permutasi membangun determinan
(\cref{ch:b2:linalg}); morfisma evaluasi $P \mapsto P(u)$
dan \omterm{def:b2:structures:ideal}{ideal} utama pada $K[X]$ melahirkan polinomial minimal serta
penguraian kernel pada \cref{ch:b2:reduction}; keterbilangan
adalah panggung tempat \cref{ch:b2:proba} bermain (peluang pada ruang
\omterm{def:b2:structures:countable}{terbilang}) sekaligus alasan mengapa topologi terus-menerus menghasilkan
himpunan padat yang \omterm{def:b2:structures:countable}{terbilang} (\cref{ch:b2:metric}). Konstruksi kuosien
$A/I$ dipakai kembali pada jilid Tahun ke-3 untuk membangun lapangan
$K[X]/(P)$ dan, dari situ, teori Galois: sifat universal yang dibuktikan
di sini dipakai di sana kata demi kata.
\end{remark}

\section{Latihan}

\begin{exercise}[$\star$]\label{exo:b2:structures:1}
Manakah di antara himpunan berikut yang \omterm{def:b2:structures:countable}{terbilang}? Himpunan semua
himpunan bagian hingga dari $\N$; himpunan \emph{semua} himpunan bagian
$\N$; $\R \setminus \Q$; himpunan polinomial berkoefisien rasional;
himpunan barisan atas $0$ dan $1$ yang akhirnya nol.
\end{exercise}

\begin{exercise}[$\star$]\label{exo:b2:structures:2}
Di $\mathfrak{S}_7$, misalkan $\sigma = (1\,4\,2\,6)(3\,5)$ dan $\tau =
(2\,3\,7)$. Hitung $\sigma\tau$ dan $\tau\sigma$ dalam bentuk \omterm{def:b2:structures:sn}{siklus}
saling lepas, \omterm{def:b2:structures:generated}{orde} dan tanda keempat permutasi tadi, serta
$\sigma^{2026}$.
\end{exercise}

\begin{exercise}[$\star$]\label{exo:b2:structures:3}
Bangun injeksi yang gamblang untuk menunjukkan bahwa $\mathcal{P}(\N)$,
$\intcc{0}{1}$ dan himpunan barisan biner $\{0,1\}^{\N}$ \omterm{def:b2:structures:countable}{ekuipoten}
sepasang demi sepasang \emph{(uraian biner dua arah;
Cantor--Bernstein menyerap kerepotan penyajian ganda)}.
\end{exercise}

\begin{exercise}[$\star$]\label{exo:b2:structures:4}
Misalkan $G$ sebuah grup dan $a, b \in G$ dua unsur komutatif yang
\omterm{def:b2:structures:generated}{ordenya} hingga dan saling prima, yakni $m$ dan $n$. Buktikan bahwa
$\operatorname{ord}(ab) = mn$. Tunjukkan lewat sebuah contoh di
$\mathfrak{S}_3$ bahwa kekomutatifan itu penting.
\end{exercise}

\begin{exercise}[$\star\star$]\label{exo:b2:structures:5}
Misalkan $G$ grup hingga berorde genap. Buktikan bahwa $G$ memuat
sebuah unsur berorde $2$. \emph{(Pasangkan tiap unsur dengan inversnya;
cacah yang berpasangan dengan dirinya sendiri.)}
\end{exercise}

\begin{exercise}[$\star\star$]\label{exo:b2:structures:6}
Buktikan bahwa $A_n$ ($n \geq 3$) \omterm{def:b2:structures:generated}{dibangun} oleh \omterm{def:b2:structures:sn}{siklus} berpanjang $3$.
\emph{(Hasil kali dua \omterm{def:b2:structures:sn}{transposisi} adalah \omterm{def:b2:structures:sn}{siklus} berpanjang $3$ atau
hasil kali dua \omterm{def:b2:structures:sn}{siklus} berpanjang $3$.)}
\end{exercise}

\begin{exercise}[$\star\star$]\label{exo:b2:structures:7}
Tentukan semua morfisma grup: dari $(\Q, +)$ ke $(\Z, +)$; dari
$(\Z/n\Z, +)$ ke $(\Z/m\Z, +)$ \emph{(cacah banyaknya:
$\gcd(m,n)$)}; dan dari $(\Q, +)$ ke $(\Q_+^*, \times)$.
\end{exercise}

\begin{exercise}[$\star\star$]\label{exo:b2:structures:8}
Dengan teorema sisa Cina, hitung $\varphi(360)$, tentukan semua
$x$ dengan $x \equiv 3 \pmod 8$, $x \equiv 5 \pmod 9$ dan $x
\equiv 2 \pmod 5$, lalu hitung dua angka terakhir $3^{2026}$
\emph{(Euler modulo $100$; awas: kerjakan modulo $4$ dan modulo $25$)}.
\end{exercise}

\begin{exercise}[$\star\star\star$]\label{exo:b2:structures:9}
Buktikan bahwa daerah integral yang hingga adalah lapangan. Turunkan
bahwa $\Z/n\Z$ adalah lapangan bila dan hanya bila $n$ prima (sekali lagi).
\end{exercise}

\begin{exercise}[$\star\star\star$]\label{exo:b2:structures:10}
(Sebuah klasik) Misalkan $K$ sebuah lapangan dan $G$ subgrup
\emph{hingga} dari $(K^*, \times)$. Buktikan bahwa $G$ \omterm{def:b2:structures:generated}{siklik}.
\emph{Petunjuk: ambil $m$ \omterm{def:b2:structures:generated}{orde} terbesar di antara unsur $G$;
tunjukkan \omterm{def:b2:structures:generated}{orde} setiap unsur membagi $m$ (dengan
\cref{exo:b2:structures:4} pada bagian saling prima yang sesuai),
sehingga seluruh $G$ memenuhi $x^m = 1$; lalu cacah akar $X^m - 1$.}
Khususnya $(\Z/p\Z)^*$ adalah \omterm{def:b2:structures:generated}{siklik}.
\end{exercise}

\begin{exercise}[$\star\star\star$]\label{exo:b2:structures:11}
Buktikan bahwa grup $(\Q, +)$ tidak \omterm{def:b2:structures:generated}{siklik}, dan yang lebih buruk lagi:
ia bahkan tidak \omterm{def:b2:structures:generated}{dibangun} oleh berhingga unsur. Buktikan sebaliknya
bahwa setiap subgrup $(\Q, +)$ yang \omterm{def:b2:structures:generated}{dibangun} oleh berhingga unsur adalah \omterm{def:b2:structures:generated}{siklik}.
\end{exercise}

\begin{exercise}[$\star\star$]\label{exo:b2:structures:12}
(Kriteria Dedekind) Buktikan bahwa setiap himpunan tak hingga memuat
himpunan bagian yang \omterm{def:b2:structures:countable}{terbilang}, lalu turunkan bahwa suatu himpunan $E$
tak hingga bila dan hanya bila ia \omterm{def:b2:structures:countable}{ekuipoten} dengan salah satu himpunan
bagian sejatinya. \emph{(Untuk arah langsungnya, geser sebuah himpunan
bagian \omterm{def:b2:structures:countable}{terbilang} sejauh satu langkah; untuk arah sebaliknya, ingat kembali asas sarang merpati.)}
\end{exercise}

\section{Soal: Teka-Teki Lima Belas}

Teka-teki lima belas adalah papan $4 \times 4$ yang memuat lima belas
ubin geser bernomor $1$ sampai $15$ dan satu sel kosong; satu langkah
menggeser salah satu ubin yang bertetangga dengan sel kosong ke sel itu.
Sekitar tahun 1890 Sam Loyd memasyhurkan teka-teki ini dengan menawarkan
\$1000 kepada siapa pun yang sanggup menukar ubin $14$ dan $15$ sambil
mengembalikan setiap ubin lain ke tempatnya. Tak seorang pun pernah
menagihnya, dan soal akhir pekan ini membuktikan kedua paruh sebabnya:
tanda permutasi pada \cref{thm:b2:structures:signature} melarang
penukaran Loyd, dan --- ini paruh yang lebih sulit sekaligus konstruktif
--- \emph{segala hal} yang diizinkan tanda permutasi memang benar-benar
terselesaikan. Pernyataan lengkapnya adalah teorema Johnson--Story (1879).

\begin{omfigure}
\begin{tikzpicture}[scale=0.72]
  \draw[omDef, thick] (0,0) grid (4,4);
  \draw[omDef, very thick] (0,0) rectangle (4,4);
  \foreach \k in {1,...,15} {
    \pgfmathtruncatemacro\row{int((\k-1)/4)}
    \pgfmathtruncatemacro\col{int(mod(\k-1,4))}
    \node at (\col+0.5, 3.5-\row) {$\k$};
  }
  \node[below, font=\small] at (2,-0.2) {tersusun};
\end{tikzpicture}
\qquad
\begin{tikzpicture}[scale=0.72]
  \draw[omDef, thick] (0,0) grid (4,4);
  \draw[omDef, very thick] (0,0) rectangle (4,4);
  \foreach \k in {1,...,13} {
    \pgfmathtruncatemacro\row{int((\k-1)/4)}
    \pgfmathtruncatemacro\col{int(mod(\k-1,4))}
    \node at (\col+0.5, 3.5-\row) {$\k$};
  }
  \node[omThm] at (1.5,0.5) {$15$};
  \node[omThm] at (2.5,0.5) {$14$};
  \node[below, font=\small] at (2,-0.2) {hadiah Loyd};
\end{tikzpicture}

{\small Konfigurasi yang tersusun dan konfigurasi $14$--$15$ milik Sam
Loyd. Pertanyaan senilai \$1000: dapatkah geseran yang sah mengubah
papan kanan menjadi papan kiri?}
\end{omfigure}

\begin{problem}[{Soal akhir pekan --- teorema keterselesaian
Johnson--Story}]
\label{pb:b2:structures:1}
Nomori selnya $1$ sampai $16$ menurut urutan baca (kiri ke kanan, atas
ke bawah), sehingga sel $k$ terletak pada baris $i$ dan kolom $j$ dengan
$k = 4(i - 1) + j$. Sel $16$ (kanan bawah) adalah \emph{rumah}
sel kosong; sel kosong kita perlakukan sebagai ubin keenam belas,
ditulis $b$ dan disamakan dengan bilangan $16$. Sebuah
\emph{konfigurasi} adalah bijeksi $\sigma \colon \intint1{16}
\to \intint1{16}$, sel $\mapsto$ isinya; konfigurasi yang
\emph{tersusun} adalah $\sigma = \mathrm{id}$. Di sepanjang soal ini,
$\varepsilon$ menyatakan tanda permutasi pada
\cref{thm:b2:structures:signature}, dan dua sel disebut
\emph{bertetangga} apabila keduanya berbagi satu rusuk papan.

\medskip
\noindent\textbf{Bagian I --- Konfigurasi, langkah, tanda.}
\begin{enumerate}
  \item Berilah alasan bahwa konfigurasi persis sama dengan unsur
        $\mathfrak{S}_{16}$, sehingga banyaknya $16! =
        20\,922\,789\,888\,000$, dan bahwa banyaknya langkah sah dari
        sebuah konfigurasi adalah $2$, $3$ atau $4$,
        bergantung pada apakah sel kosong berada di pojok, di
        tepi, atau di bagian dalam.
  \item Misalkan $\sigma$ sebuah konfigurasi, $p = \sigma^{-1}(16)$ sel
        tempat kekosongan berada, dan $c$ sel yang bertetangga dengan
        $p$. Tunjukkan bahwa menggeser ubin di $c$ ke $p$ menghasilkan
        konfigurasi $\sigma' = \sigma \circ \tau$ dengan $\tau =
        (p\ c)$, lalu turunkan bahwa setiap langkah membalik tandanya:
        $\varepsilon(\sigma') = -\varepsilon(\sigma)$.
  \item Warnai papan bak papan catur: $\chi(k) = (-1)^{i+j}$ untuk sel
        $k$ pada baris $i$, kolom $j$. Tunjukkan bahwa setiap langkah
        membalik $\chi(\text{sel tempat kekosongan})$, lalu turunkan
        bahwa rangkaian langkah yang mengembalikan kekosongan ke sel
        awalnya pasti berpanjang genap.
  \item Tunjukkan bahwa
        \[
        I(\sigma) = \varepsilon(\sigma)\,
        \chi\bigl(\sigma^{-1}(16)\bigr)
        \]
        tidak berubah oleh langkah sah mana pun, lalu hitung
        $I(\mathrm{id})$.
\end{enumerate}

\medskip
\noindent\textbf{Bagian II --- Hadiah Loyd: invarian sedang
bekerja.}
\begin{enumerate}[resume]
  \item Konfigurasi Loyd $\sigma_L$ sama dengan konfigurasi tersusun
        kecuali bahwa sel $14$ dan $15$ memuat ubin $15$ dan $14$.
        Hitung $I(\sigma_L)$ lalu simpulkan bahwa tidak ada rangkaian
        langkah yang menghubungkan $\sigma_L$ dengan konfigurasi
        tersusun: \$1000 milik Loyd tidak pernah terancam.
  \item Tunjukkan bahwa tepat separuh dari semua konfigurasi memenuhi
        $I = +1$: $\abs{\{\sigma : I(\sigma) = +1\}} = 16!/2$.
        \emph{(Untuk sel kosong yang tetap, pasangkan konfigurasi
        dengan menyusunnya bersama satu \omterm{def:b2:structures:sn}{transposisi} tetap atas dua sel
        yang lain.)}
  \item Tunjukkan bahwa setiap langkah dapat dibatalkan oleh sebuah
        langkah sah, bahwa ``$\sigma'$ terjangkau dari $\sigma$ lewat
        langkah sah'' merupakan relasi ekuivalensi, dan bahwa kelas $R$
        milik konfigurasi tersusun memenuhi $R \subseteq \{I = +1\}$.
        Simpulkan bahwa kelasnya sedikitnya ada dua.
  \item Andaikan kekosongan berada di rumahnya: $\sigma(16) = 16$.
        Tunjukkan bahwa $I(\sigma) = \varepsilon(\rho)$ dengan $\rho \in
        \mathfrak{S}_{15}$ adalah pembatasan $\sigma$ pada
        sel $1, \dots, 15$, dan bahwa sembarang konfigurasi dapat
        dibawa oleh langkah sah ke konfigurasi yang kekosongannya di
        rumah. Simpulkan: untuk membuktikan $R = \{I = +1\}$ cukuplah
        mewujudkan setiap permutasi \emph{genap} atas kelima belas sel
        bukan-rumah lewat rangkaian langkah yang berawal dan berakhir
        dengan kekosongan di rumah.
\end{enumerate}

\medskip
\noindent\textbf{Bagian III --- Perjalanan kekosongan dan grup
program.} Sebuah \emph{program} adalah rangkaian langkah sah yang
hingga, dimulai dari konfigurasi yang kekosongannya di rumah, dan
konfigurasi akhirnya pun berkekosongan di rumah. \emph{Efeknya} adalah
permutasi $\pi$ atas sel yang ditetapkan oleh: isi sel $x$ berakhir di
sel $\pi(x)$.
\begin{enumerate}[resume]
  \item Tunjukkan bahwa program yang dijalankan dari $\sigma$ berakhir
        di $\sigma \circ \pi^{-1}$; bahwa menjalankan dua program
        berturut-turut menyusun efek keduanya; dan bahwa himpunan $H$
        semua efek adalah subgrup $\mathfrak{S}_{15}$ (permutasi
        atas sel $1, \dots, 15$) yang termuat di grup
        alternating $A_{15}$.
  \item (Perjalanan dasar) Dari kekosongan di rumah, geserlah
        kekosongan mengelilingi blok $2 \times 2$ di kanan bawah: sel
        $16 \to 12 \to 11 \to 15 \to 16$. Tunjukkan efeknya adalah
        \omterm{def:b2:structures:sn}{siklus} berpanjang $3$, yakni $(11\ 12\ 15)$, dan bahwa perjalanan
        sebaliknya memberi
        $(11\ 15\ 12)$. Keduanya berada di $H$.
  \item (Perjalanan agung) Periksalah bahwa
        \[
        16 \to 15 \to 14 \to 13 \to 9 \to 5 \to 1 \to 2 \to 3
        \to 4 \to 8 \to 7 \to 6 \to 10 \to 11 \to 12 \to 16
        \]
        adalah jalan tertutup yang melewati keenam belas sel (hanya
        lewat langkah bertetangga), dan bahwa efeknya adalah \omterm{def:b2:structures:sn}{siklus}
        berpanjang $15$
        \[
        \zeta = (15\ 12\ 11\ 10\ 6\ 7\ 8\ 4\ 3\ 2\ 1\ 5\ 9\ 13\
        14) .
        \]
        Dengan menulis $x_0 = 15$, $x_1 = 12$, $x_2 = 11$, \dots, $x_{14}
        = 14$ menurut urutan siklusnya, periksalah bahwa perjalanan
        dasar yang dibalik pada pertanyaan 10 persis sama dengan $(x_0\ x_1\
        x_2)$.
  \item Buktikan rumus konjugasi pada sembarang $\mathfrak{S}_n$: untuk
        sebuah permutasi $g$ dan sebuah \omterm{def:b2:structures:sn}{siklus} berpanjang $3$,
        \[
        g\,(a\ b\ c)\,g^{-1} = \bigl(g(a)\ g(b)\ g(c)\bigr),
        \]
        lalu perhatikan bahwa $H$, karena ia grup, tertutup terhadap
        konjugasi oleh unsurnya sendiri.
  \item Turunkan bahwa $H$ memuat kelima belas \omterm{def:b2:structures:sn}{siklus} \emph{berurutan}
        berpanjang $3$ dari perjalanan agung:
        \[
        s_t = (x_t\ x_{t+1}\ x_{t+2}) \qquad (t \in \Z/15\Z,
        \text{ indeks modulo } 15).
        \]
\end{enumerate}

\medskip
\noindent\textbf{Bagian IV --- Membangun grup alternating.}
\begin{enumerate}[resume]
  \item (Lema A) Misalkan $s$ dan $t$ dua \omterm{def:b2:structures:sn}{siklus} berpanjang $3$ yang
        penyangganya berbagi tepat dua titik, katakanlah penyangga
        $\{a, b, c\}$ dan $\{b, c, d\}$. Tunjukkan bahwa, setelah $s$
        atau $t$ diganti dengan inversnya bila perlu (yang tidak
        mengubah subgrup yang \omterm{def:b2:structures:generated}{dibangun}), hasil kali $st$ adalah
        \omterm{def:b2:structures:sn}{transposisi} ganda; tunjukkan bahwa $A_4$ tidak punya subgrup
        berorde $6$ \emph{(subgrup berindeks $2$ memuat setiap
        kuadrat; cacah \omterm{def:b2:structures:sn}{siklus} berpanjang $3$ di antara kuadrat itu)};
        lalu simpulkan bahwa $\langle s, t\rangle$ adalah seluruh grup
        alternating atas keempat huruf $\{a, b, c, d\}$.
  \item (Lema B) Misalkan $X$ himpunan berisi $k \geq 4$ huruf, $w
        \notin X$, dan misalkan $G$ subgrup suatu
        $\mathfrak{S}_n$ yang memuat setiap permutasi genap atas $X$
        dan satu \omterm{def:b2:structures:sn}{siklus} berpanjang $3$, yaitu $(u\ v\ w)$, dengan $u, v \in X$.
        Tunjukkan bahwa untuk setiap $a, b \in X$ yang berbeda ada
        permutasi \emph{genap} $g$ atas $X$ dengan $g(u) = a$, $g(v) =
        b$, lalu turunkan $(a\ b\ w) \in G$.
  \item Turunkan bahwa grup $G$ pada Lema B memuat setiap permutasi
        genap atas $X \cup \{w\}$ \emph{(pakai
        \cref{exo:b2:structures:6}: \omterm{def:b2:structures:sn}{siklus} berpanjang $3$ membangunnya)}.
        Lalu, dengan merangkai Lema A dan Lema B sepanjang \omterm{def:b2:structures:sn}{siklus}
        berpanjang $3$ yang berurutan, yaitu $s_0, s_1, \dots, s_{12}$ pada
        pertanyaan 13,
        buktikan bahwa $\langle s_0, \dots, s_{12}\rangle = A_{15}$.
  \item Simpulkan bahwa $H = A_{15}$: \emph{setiap penataan ulang yang
        genap atas kelima belas ubin dapat dicapai oleh sebuah program},
        dan $H$ punya $15!/2 = 653\,837\,184\,000$ unsur.
  \item (Teorema Johnson--Story, 1879) Rakitlah pertanyaan 6,
        7, 8 dan 17: konfigurasi yang terjangkau dari konfigurasi
        tersusun \emph{persis} sebanyak $16!/2 =
        10\,461\,394\,944\,000$ konfigurasi dengan $I = +1$;
        dan keterjangkauan hanya punya \emph{dua} kelas, yakni kelas
        konfigurasi tersusun dan kelas $\sigma_L$ milik
        Loyd. \emph{(Untuk butir kedua, tukarlah nama ubin
        $14$ dan $15$: tunjukkan $\sigma \mapsto (14\ 15) \circ
        \sigma$ memetakan rangkaian langkah ke rangkaian langkah dan
        menukar $\{I = +1\}$ dengan $\{I = -1\}$.)}
\end{enumerate}

\medskip
\noindent\textbf{Bagian V --- Kriteria, ragam, dan pandangan
dari atas.}
\begin{enumerate}[resume]
  \item (Kriteria praktis) Bacalah kelima belas ubin itu menurut urutan
        baca selnya, lewati kekosongannya, dan misalkan
        $N$ banyaknya inversi pada deretan itu; misalkan $r$ nomor
        baris kekosongan dihitung dari \emph{bawah}. Tunjukkan
        bahwa $I(\sigma) = (-1)^{N + r + 1}$, sehingga $\sigma$
        terselesaikan bila dan hanya bila $N + r$ ganjil.
  \item (Aksi grup) Sebuah \emph{aksi} grup $G$ pada himpunan
        $X$ adalah pemetaan $G \times X \to X$, $(g, x) \mapsto g \cdot
        x$, dengan $e \cdot x = x$ dan $g \cdot (h \cdot x) =
        (gh) \cdot x$; \emph{orbit} $x$ adalah $G \cdot x$, dan
        aksinya disebut \emph{bebas} apabila $g \cdot x = x$ memaksa $g =
        e$. Tunjukkan bahwa $h \cdot \sigma = \sigma \circ h^{-1}$
        mendefinisikan aksi bebas $H$ pada himpunan konfigurasi yang
        kekosongannya di rumah, bahwa orbitnya persis kelas
        keterjangkauan bersama lewat program, lalu peroleh kembali dari
        pencacahan orbit bahwa konfigurasi itu terbelah menjadi tepat
        $15!\,/\,\abs H = 2$ kelas.
  \item (Halangan pada papan $3 \times 3$) Tunjukkan bahwa papan
        $3 \times 3$ \emph{tidak} punya jalan tertutup yang menyinggahi
        setiap sel tepat sekali: siasat perjalanan agung pada Bagian III
        gagal untuk teka-teki delapan. \emph{(Warnai kesembilan
        selnya bak papan catur.)}
  \item (Perbaikannya) Pada papan $3 \times 3$ dengan sel $1$ sampai
        $9$ menurut urutan baca dan rumah di $9$: hitunglah efek
        perjalanan keliling $9 \to 8 \to 7 \to 4 \to 1 \to 2 \to 3
        \to 6 \to 9$ (sebuah \omterm{def:b2:structures:sn}{siklus} berpanjang $7$, yaitu $\zeta'$, yang
        membiarkan pusat $5$ tetap) dan efek perjalanan pojok $9 \to 6
        \to 5 \to 8 \to 9$ (\omterm{def:b2:structures:sn}{siklus} berpanjang $3$ yang melewati pusat).
        Dengan mengonjugasikan yang terakhir oleh pangkat-pangkat
        $\zeta'$ lalu merangkai Lema A dan Lema B, buktikan
        bahwa grup program teka-teki delapan adalah seluruh $A_8$,
        sehingga tepat $9!/2 = 181\,440$ dari $9! =
        362\,880$ konfigurasinya terselesaikan.
  \item (Papan yang miskin) Kini ambil papan berupa satu \omterm{def:b2:structures:sn}{siklus}
        beranggotakan $n \geq 4$ sel yang mengusung $n - 1$ ubin.
        Tunjukkan bahwa urutan siklis ubinnya tidak berubah, bahwa tiap
        kelas keterjangkauan punya tepat $n(n - 1)$ konfigurasi
        \emph{(kelasnya adalah orbit \omterm{def:b2:structures:generated}{grup siklik} berorde
        $\operatorname{lcm}(n, n-1) = n(n-1)$)}, dan bahwa kelasnya ada
        $(n - 2)!$ --- untuk $n \geq 5$ jauh lebih banyak daripada
        $2$: pada papan yang tipis invarian paritas nyaris tak
        menangkap apa pun, dan geometrilah yang berkuasa.
  \item Dua vonis menurut kriteria pertanyaan 19: papan yang
        seluruhnya terbalik (ubin $15, 14, \dots, 1$ pada sel $1$
        sampai $15$, kekosongan di rumah) dan papan yang kekosongannya
        di sel $1$ lalu disusul ubin $15, 14, \dots, 1$ pada sel $2$
        sampai $16$. Manakah yang terselesaikan?
  \item (Rangkuman) Buktinya bertumpu pada dua pilar yang saling bebas:
        sebuah \emph{invarian} ($I$, yang \omterm{def:b2:structures:generated}{dibangun} dari morfisma tanda
        permutasi) yang menunjukkan paling banyak separuh konfigurasi
        terjangkau, dan sebuah teorema \emph{pembangunan gamblang}
        ($H = A_{15}$) yang menunjukkan sedikitnya separuh terjangkau.
        Dengan satu kalimat untuk masing-masing, sebutkan di mana
        berikut ini masuk: sifat morfisma
        $\varepsilon$; teorema Lagrange; pembangunan $A_n$
        oleh \omterm{def:b2:structures:sn}{siklus} berpanjang $3$; dan konjugasi. Nyatakan asas
        umumnya dalam satu baris.

\end{enumerate}
\end{problem}
